%%%%%%%%%%%%%%%%%%%%%%%%%%%%%%%%%%%%%%%%%%%%%%%%%%%%%%%%%%%%%%%%%%%%%%%%%%%%%%%
% Präsi für Kolloquium
%%%%%%%%%%%%%%%%%%%%%%%%%%%%%%%%%%%%%%%%%%%%%%%%%%%%%%%%%%%%%%%%%%%%%%%%%%%%%%%

\documentclass[aspectratio=169]{beamer}

\usepackage[ngerman]{babel}
\usepackage{fontspec}
\usepackage{microtype}
\usepackage{selnolig}
\usepackage{csquotes}
\usepackage{amsmath, amssymb}
\usepackage{bm}
\usepackage{soul}
\usepackage{booktabs}
\usepackage{mwe}
\usepackage[style=ieee, backend=biber]{biblatex}
\usepackage{tikz}
\usetikzlibrary{
	positioning,
	calc,
	decorations.shapes,
	decorations.pathmorphing,
	decorations.pathreplacing,
}
\usepackage{listings}


\graphicspath{{img/}}


\addbibresource{refs.bib}


\hypersetup{
	pdfauthor={André Ulrich},
	pdftitle={Kolloquium zur Masterarbeit},
}


\lstset{
	basicstyle=\ttfamily,
	morekeywords={
		{\verb|\begin|},
		\textbackslash end,
		\textbackslash documentclass,
	},
	keywordstyle=\color{blue},
	frame=single,
}


\newcommand*{\param}[1]{$\langle$\textit{\ttfamily #1}$\rangle$}
\newcommand*{\tb}{\textbackslash}

\makeatletter
    \newenvironment{withoutheadline}{
        \setbeamertemplate{headline}[default]
        \def\beamer@entrycode{\vspace*{-\headheight}}
    }{}
\makeatother


\pgfdeclarelayer{background}
\pgfsetlayers{background, main}


\usetheme{frankfurt}
\setbeamertemplate{footline}[frame number]
\setbeamertemplate{navigation symbols}{}


\begin{document}
\author{André Ulrich}
\title{Algorithmus zur optimierten Ladung von Elektrofahrzeugen in einem
gemeinsamen Netzzweig}
\subtitle{Master Kolloquium}
\date{
	20. September 2022\bigskip
	
	\begin{tabular}{@{}ll@{}}
	Betreuer: & Prof.\,Dr. Eberhard Waffenschmidt \\
	Ko-Referent: & Prof.\,Dr. Ingo Stadler \\
	\end{tabular}
}
\titlegraphic{
	\pgfputat{\pgfxy(6,0.4)}{
		\includegraphics[scale=0.05]{TH_Koeln_Logo.pdf}
	}
}


\setbeamertemplate{footline}[default]
\begin{withoutheadline}
\begin{frame}[noframenumbering]
	\maketitle
\end{frame}
\end{withoutheadline}


%\setbeamertemplate{footline}[frame number]
\begin{withoutheadline}
\begin{frame}[noframenumbering]{Inhalt}
	\tableofcontents
\end{frame}
\end{withoutheadline}


\section{Hintergrund und Aufgabenstellung}
\setbeamertemplate{footline}[frame number]

\begin{frame}{Hintergrund}
	\begin{itemize}
		\item PROGRESSUS-Projekt: weniger Netzüberlastung, mehr BEVs laden
		\cite{noauthor_progressus_nodate}
		\item Entwicklung von Optimierungsalgorithmen zur koordinierten Ladung 
		von BEVs
		\item Bestehende Lösungsansätze:
		\begin{itemize}
			\item GridTopologyEstimation (GTE) zum Bestimmen der Netztopologie
			\item GridStateEstimation (GSE) zum Bestimmen des Netzzustands
		\end{itemize}
		\item \so{ABER:} Bisher nur eingeschränkte Ansätze zur Optimierung
		wie \cite{iserloh2021} oder \cite{koese2022} vorhanden
	\end{itemize}
\end{frame}


\begin{frame}{Hintergrund}
	\begin{figure}
		\centering
		%%%%%%%%%%%%%%%%%%%%%%%%%%%%%%%%%%%%%%%%%%%%%%%%%%%%%%%%%%%%%%%%%%%%%%%%%%%%%%%
% Bild zum Überblick der aktuellen Aufgaben im PROGRESSUS-Projekt
%%%%%%%%%%%%%%%%%%%%%%%%%%%%%%%%%%%%%%%%%%%%%%%%%%%%%%%%%%%%%%%%%%%%%%%%%%%%%%%

\begin{tikzpicture}[
	every node/.style={
		font=\footnotesize
	},
	block/.style={
		draw,
		rectangle,
		align=center,
		minimum width=1cm,
		execute at begin node=\strut
	}
]

	\node[block, rounded corners=3pt](SMG){\strut SMG};
	
	\node[block, right=2.5 of SMG](GTE){GTE};
	
	\node[block, above=2 of SMG, rounded corners=3pt](hardware){Hardware};
	
	\node[block, above=2 of GTE](software){Software};
	
	\node[block, right=1.5 of GTE](GSE){GSE};
	
	\node[block, fill=orange!50,  right=1.5 of GSE , minimum height=3.5cm,
	minimum width=1.5cm](GLO){\strut GLO};
	
	\node[block, right=1.5cm of GLO, minimum height=2.5cm,
	minimum width=1.5cm](EMO){EMO};
	
	\node[block, above=0.8cm of EMO, fill opacity=0, draw opacity=0](GLS){GLS};
	
	
	% Pfeil von SMG zu GSE
	\draw[-stealth] (SMG.270) -- +(0,-0.4) -| (GSE.240) node[pos=0.25, below]
	{$I$, $U$, $\cos\varphi $};
	
	\draw[-stealth] (SMG.east) -- (GTE.west) node[pos=0.5, below]
	{$I$, $U$, $\cos\varphi $};
	
	\draw[-stealth] (GTE.east) -- (GSE.west) node[pos=0.5, below]
	{$\bm{Y}$};
	
	\draw[dashed, -stealth] (GTE.north) |- (GLO.135) node[pos=0.75, above]
	{$\bm{Y}$};
	
	\draw[dashed, -stealth] (GSE.east) -- (GLO.west) node[pos=0.5, below]
	{$I^\mathsf{HH}$};
	
	\draw[dashed, -stealth] (GLO.225) -| (GSE.300)
	node[pos=0.25, below]{$I^\mathsf{BEV}$};
	
%	\draw[-stealth] (EMO.135) -- (GLO.45) node[pos=0.5, above]
%	{$U^\mathrm{Knoten}$};

	\draw[-stealth] (EMO.135) -- (EMO.135 -| GLO.east)
	node[pos=0.5, above]{$U^\mathsf{Knoten}$};
	
	\draw[-stealth] (GLO.315) -- (EMO.225) node[pos=0.5, above]
	{$I^\mathsf{BEV}$};
	
	\draw[-stealth] (GLS.west) -| (GLO.north) node[pos=0.25, above]
	{Szenarien};
	
	\begin{pgfonlayer}{background}
		\path(GLO.south west)+(-0.2,-0.8) node (helper-a){};
		\path(GLO.north -| EMO.east)+(0.2,1.3) node (helper-b){};
		\path[fill=yellow!20, rounded corners, draw=black!50, dashed]
		(helper-a) rectangle (helper-b);
	\end{pgfonlayer}
	
	\node at ($($(helper-a) !0.5! (helper-b)$) - (0,2.5)$){Teil dieser Arbeit};

\end{tikzpicture}
	\end{figure}
\end{frame}


\begin{frame}{Aufgabenstellung}
	\begin{enumerate}
		\item Entwicklung eines Optimierungsalgorithmus zur optimalen 
		Koordination der Ladevorgänge von BEVs
		\item Verknüpfung von Optimierung und Simulationsumgebung
		\item Ermöglichen eines parallelen Betriebs von Optimierung und
		Simulationsumgebung
	\end{enumerate}
\end{frame}


\begin{frame}{Aufgabenstellung}
	Der Optimierungsalgorithmus soll:
	\begin{itemize}
		\item Zunächst nur einzelne Netzstrahlen betrachten
		\item Netzauslastung maximieren, um möglichst viele BEVs zu laden
		\item Gleichzeitig jedoch Netzüberlastung verhindern
		\item Nach Möglichkeit Diskriminierung beim Laden ausschließen
		\item Automatisch skalieren mit dem betrachteten Netzstrahl
		\item Auch auf einem Raspberry Pi in \enquote{vertretbarer} Zeit 
		ausführbar sein
	\end{itemize}
	
	\onslide<2>{
		\begin{block}{also}
			Einschränkungen bisheriger Arbeiten überwinden und 
			\enquote{Gratwanderung vollführen}.
		\end{block}
	}
\end{frame}

\section{Theorie}

\begin{frame}{Prinzip des Optimierungsalgorithmus}
	\begin{figure}
		\centering
		%%%%%%%%%%%%%%%%%%%%%%%%%%%%%%%%%%%%%%%%%%%%%%%%%%%%%%%%%%%%%%%%%%%%%%%%%%%%%%%
% Bild für die Input-/Output-Parameter
%%%%%%%%%%%%%%%%%%%%%%%%%%%%%%%%%%%%%%%%%%%%%%%%%%%%%%%%%%%%%%%%%%%%%%%%%%%%%%%

\begin{tikzpicture}[
		block/.style={
			font=\tiny,
			draw,
			rectangle,
			minimum width={3.5cm}
			}
		]
		%\node[block](block0){\strut Haushaltslastprofile};
		
		\node[block](block1){\strut Transformatorleistung};		
		
		\node[block,
			below=0.2 of block1](block2){\strut Anzahl der Knoten};
			
		\node[block,
			below=0.2 of block2](block3){\strut Anzahl Ladestationen};
			
%		\node[block,
%			below=0.2 of block3](block4){\strut Leitungsimpedanzen};
			
		\node(gridtopology)[below=0.1 of block3]{\tiny\strut 
		Netztopologie};
			
		\node[block,
			below=0.3 of gridtopology](block5){\strut Start/Ziel SOC};
			
		\node[block,
			below=0.2 of block5](block6){\strut Start/Ziel Zeit};
			
		\node(customer)[below=0.1 of block6]{\tiny\strut Kundenwünsche};
			
		\node[block,
			below=0.3 of customer](block7){\strut Zeitliche Auflösung};
			
		\node[block,
			below=0.2 of block7](block8){\strut Weite des Horizonts};
			
		\node(timespecs)[below=0.1 of block8]{\tiny\strut Zeitliche
		Betrachtung};
			
		\node[block,
			right=2 of block5, align=center,
			minimum height=5cm](GLO){GridLineOptimizer\\ (GLO)};
			
		\node[block,
			right=of GLO, align=left](res){Zeitreihen für \\ Ladestrom und 
			SOC};
			
		\node[block,
			below=0.5 of GLO](block0){\strut Haushaltslastprofile};
			
		
		\draw[-stealth](block0.north) -- (GLO);
		\draw[-stealth](block1.east) -- (GLO);
		\draw[-stealth](block2.east) -- (GLO);
		\draw[-stealth](block3.east) -- (GLO);
		%\draw[-stealth](block4.east) -- (GLO);
		\draw[-stealth](block5.east) -- (GLO);
		\draw[-stealth](block6.east) -- (GLO);
		\draw[-stealth](block7.east) -- (GLO);
		\draw[-stealth](block8.east) -- (GLO);
		\draw[-stealth](GLO) -- (res);
		
		
		\begin{pgfonlayer}{background}
			\path(block1.west)+(-0.1, 0.4) node (a) {};
			\path(block3.east)+(0.1, -0.9) node (b) {};
			\path[fill=yellow!20, rounded corners, draw=black!50, dashed]
			(a) rectangle (b);
			
			\path(block5.west)+(-0.1, 0.4) node (a) {};
			\path(block6.east)+(0.1, -0.9) node (b) {};
			\path[fill=red!20, rounded corners, draw=black!50, dashed]
			(a) rectangle (b);
			
			\path(block7.west)+(-0.1, 0.4) node (a) {};
			\path(block8.east)+(0.1, -0.9) node (b) {};
			\path[fill=green!20, rounded corners, draw=black!50, dashed]
			(a) rectangle (b);
		\end{pgfonlayer}
	\end{tikzpicture}
	\end{figure}
\end{frame}


\begin{frame}{Getroffene Annahmen}
	\begin{itemize}
		\item Nur Verbraucher im Netzstrahl
		\item Nur ohmschen Anteil der Leistungsimpedanzen berücksichtigen
		\item Hausanschlussleitungen vernachlässigen
	\end{itemize}
	
	\onslide<2>{
		\begin{block}{Vereinfachte Darstellung}
			\begin{itemize}
				\item Hinreichende Beschreibung des Netzstrahls durch
				eine Reihe von Widerständen
				\item Betrachtung der Spannung am letzten Knoten ausreichend
			\end{itemize}
		\end{block}
	}
\end{frame}


\begin{frame}{Generischer Netzstrahl}
	Aus den getroffenen Vereinfachungen resultierender, generischer Netzstrahl:
	\vspace{1cm}
	
	\begin{figure}
		%%%%%%%%%%%%%%%%%%%%%%%%%%%%%%%%%%%%%%%%%%%%%%%%%%%%%%%%%%%%%%%%%%%%%%%%%%%%%%%
% Bild für die Netzstrang
%%%%%%%%%%%%%%%%%%%%%%%%%%%%%%%%%%%%%%%%%%%%%%%%%%%%%%%%%%%%%%%%%%%%%%%%%%%%%%%

\begin{tikzpicture}[node/.style={fontsize=\scriptsize}]
	% Koordinaten der wichtigen Punkte
	\coordinate (TrafoHi) at (0,0);
	\coordinate (TrafoLo) at (0.3,0);
	\coordinate (LineStart) at (0.7,0);
	\coordinate (LineEnd) at (5,0);
	\coordinate (SecLineStart) at (7,0);
	\coordinate (SecLineEnd) at (10,0);
	% Knoten
	\coordinate (Node1) at (2,0);
	\coordinate (Node2) at (4,0);
	\coordinate (Node3) at (8,0);
	\coordinate (Node4) at (10,0);
	% Leitungen (deren Beschriftung)
	\coordinate (FirstLineAnchor) at (1,0);
	\coordinate (SecondLineAnchor) at ($(Node1)!.5!(Node2)$);
	\coordinate (ThirdLineAnchor) at ($(Node2)!.25!(Node3)$);
	\coordinate (FourthLineAnchor) at ($(Node2)!.75!(Node3)$);
	\coordinate (FifthLineAnchor) at ($(Node3)!.5!(Node4)$);
	
	% unten der abstrakte Netzstrahl
	\draw (TrafoHi) circle(4mm);
	\draw (TrafoLo) circle(4mm);
	\draw (LineStart) -- (LineEnd);
	\draw (SecLineStart) -- (SecLineEnd);
	\draw[dashed] (LineEnd) -- (SecLineStart);
	% und alle Knoten
	\fill (Node1) circle(0.5mm) node[above]{$n_1$};
	\fill (Node2) circle(0.5mm) node[above]{$n_2$};
	\fill (Node3) circle(0.5mm) node[above]{$n_{n-1}$};
	\fill (Node4) circle(0.5mm) node[above]{$n_n$};
	% sowie alle Leitungen
	\node[below] at (FirstLineAnchor) {$l_1$};
	\node[below] at (SecondLineAnchor) {$l_2$};
	\node[below] at (ThirdLineAnchor) {$l_3$};
	\node[below] at (FourthLineAnchor) {$l_{n-1}$};
	\node[below] at (FifthLineAnchor) {$l_n$};
	
	% oben der spezielle Netzstrahl
	\draw ($(TrafoHi)+(0,3)$) circle(4mm);
	\draw ($(TrafoLo)+(0,3)$) circle(4mm);
	\draw ($(LineStart)+(0,3)$) -- ($(LineEnd)+(0,3)$);
	\draw ($(SecLineStart)+(0,3)$) -- ($(SecLineEnd)+(0,3)$);
	\draw[dashed] ($(LineEnd)+(0,3)$) -- ($(SecLineStart)+(0,3)$);
	% und alle Knoten
	\draw[ultra thick] ($(Node1)+(0,3.4)$) -- ($(Node1)+(0,2.6)$);
	\draw[ultra thick] ($(Node2)+(0,3.4)$) -- ($(Node2)+(0,2.6)$);
	\draw[ultra thick] ($(Node3)+(0,3.4)$) -- ($(Node3)+(0,2.6)$);
	\draw[ultra thick] ($(Node4)+(0,3.4)$) -- ($(Node4)+(0,2.6)$);
	% Abzweige zu den jeweiligen Lasten
	\draw[-stealth] ($(Node1)+(0,3)$) -- +(0.5,-0.5) -- +(0.5,-1) node[below]
	{$I_1$};
	\draw[-stealth] ($(Node2)+(0,3)$) -- +(0.5,-0.5) -- +(0.5,-1) node[below]
	{$I_2$};
	\draw[-stealth] ($(Node3)+(0,3)$) -- +(0.5,-0.5) -- +(0.5,-1) node[below]
	{$I_{n-1}$};
	\draw[-stealth] ($(Node4)+(0,3)$) -- +(0.5,-0.5) -- +(0.5,-1) node[below]
	{$I_n$};
	\node[above] at ($(FirstLineAnchor)+(0,3)$) {$Z_1$};
	\node[above] at ($(SecondLineAnchor)+(0,3)$) {$Z_2$};
	\node[above] at ($(ThirdLineAnchor)+(0,3)$) {$Z_3$};
	\node[above] at ($(FourthLineAnchor)+(0,3)$) {$Z_{n-1}$};
	\node[above] at ($(FifthLineAnchor)+(0,3)$) {$Z_n$};
	% Spannungen an den Knoten
	\draw[-stealth] ($(Node1)+(-0.1,2.9)$) .. controls +(-0.5,-0.5) .. +(0,-1)
	node[left]{$U_1$};
	\draw[-stealth] ($(Node2)+(-0.1,2.9)$) .. controls +(-0.5,-0.5) .. +(0,-1)
	node[left]{$U_2$};
	\draw[-stealth] ($(Node3)+(-0.1,2.9)$) .. controls +(-0.5,-0.5) .. +(0,-1)
	node[left]{$U_{n-1}$};
	\draw[-stealth] ($(Node4)+(-0.1,2.9)$) .. controls +(-0.5,-0.5) .. +(0,-1)
	node[left]{$U_n$};
\end{tikzpicture}
	\end{figure}
\end{frame}


\begin{frame}{Optimierung mit Mathematical Programming}
	\begin{columns}
		\begin{column}{.6\textwidth}
			\begin{itemize}
				\item Mathematical Programming als strukturierte 
				Darstellungsweise eines Optimierungsproblems
				\item Formulierug als Linear Programmig
				\item Lösung mithilfe von geeignetem Solver
			\end{itemize}
		\end{column}
		\begin{column}{.4\textwidth}
			\begin{align*}
				\begin{split}
		\text{max.}\quad&\sum_{i\in\mathcal{I}}c_i\cdot x_i \\
		\text{s.\,t.}\quad&\sum_{i\in\mathcal{I}}a_{j,i}\cdot x_i\leq b_j\quad
		\forall j\in\mathcal{J} \\
		& x_i\in\mathbb{R}_0^+\quad\forall i\in\mathcal{I}
				\end{split}
			\end{align*}
		\end{column}
	\end{columns}
\end{frame}


\begin{frame}{Grundbausteine eines Mathematical Programming}
	\begin{description}[Restriktionen]
		\item[Sets] zum indexieren der restlichen Bestandteile
		\item[Parameter] als festgelegte Größen (beeinflussen Optimierung)
		\item[Variablen] als zu bestimmende, variable Größen
		\item[Zielfunktion] in Abhängigkeit der Variablen (zu minimieren oder
		maximieren)
		\item[Restriktionen] zum Einschränken der möglichen Lösungsmenge
	\end{description}
	
	\onslide<2>{
		\begin{block}{Ziel:}
			Variablen so bestimmen, dass Zielfunktion maximiert/minimiert wird
		\end{block}
	}
\end{frame}


\begin{frame}{Formulierung als Linear Programming}
	\begin{Large}
		Sets:
	\end{Large}
	\begin{itemize}
		\item Menge aller Knoten: $\mathcal{N}=\{n_1, n_2,\dots,n_n\}$
		\item Menge aller Leitungen: $\mathcal{L}=\{l_1,l_2,\dots,l_n\}$
		\item Menge aller Koten mit Ladesäule und BEV:
		$\mathcal{N}_\mathsf{BEV}\subset\mathcal{N}$
		\item Menge aller Zeitschritte: $\mathcal{T}=\{t_1,t_2,\dots,t_n\}$
	\end{itemize}
%	\onslide<2>{
%		\begin{block}{Weil:}
%			Alle relevanten Größen lassen sich in den Sets indexieren
%		\end{block}
%	}
\end{frame}


\begin{frame}{Formulierung als Linear Programming}
	\begin{Large}
		Parameter:
	\end{Large}
	\begin{itemize}
		\item Leitungsimpedanzen: $Z_l\quad\forall l\in\mathcal{L}$
		\item Knotenspannungen: $U_n\quad\forall n\in\mathcal{N}$
		\item Haushaltsströme:
		$I^\mathsf{HH}_{t,n}\quad\forall t,n\in\mathcal{T}\times\mathcal{N}$
		\item Start SOCs der BEVs:
		$SOC^\mathsf{Start}_{n}\quad\forall n\in\mathcal{N}_\mathsf{BEV}$
		\item Ziel SOCs der BEVs:
		$SOC^\mathsf{Ziel}_{n}\quad\forall n\in\mathcal{N}_\mathsf{BEV}$
		\item $\dots$
	\end{itemize}
%	\onslide<2>{
%		\begin{block}{Weil:}
%			Festgelegte Größen, die das Ergebnis der Optimierung beeinflussen
%		\end{block}
%	}
\end{frame}


\begin{frame}{Formulierung als Linear Programming}
	\begin{Large}
		Variablen:
	\end{Large}
	\begin{itemize}
		\item Ladeströme der BEVs:
		$I^\mathsf{BEV}_{t,n}\quad\forall t,n\in\mathcal{T}\times
		\mathcal{N}_\mathsf{BEV}$
		\item SOCs der BEVs:
		$SOC_\mathsf{t,n}\quad\forall t,n\in\mathcal{T}\times
		\mathcal{N}_\mathsf{BEV}$
	\end{itemize}
%	\onslide<2>{
%		\begin{block}{Weil:}
%			Das sind die von der Optimierung zu bestimmenden Größen
%		\end{block}
%	}
\end{frame}


\begin{frame}{Formulierung als Linear Programming}
	\begin{Large}
		Zielfunktion:
	\end{Large}\vspace{.15cm}
	
	Maximieren der Ladeströme.
	\begin{equation*}
		\mathsf{max.}\quad \sum_{t\in\mathcal{T}}\Big(\sum_{n\in
		\mathcal{N}_\mathsf{BEV}}\!\!I^\mathsf{BEV}_{t,n}\Big)
	\end{equation*}
%	\onslide<2>{
%		\begin{block}{Weil:}
%			Das ist das Ziel der Optimierung
%		\end{block}
%	}
\end{frame}


\begin{frame}[t]{Formulierung als Linear Programming}
	\vspace*{1cm}
	
	\begin{Large}
		Restriktionen:
	\end{Large}\vspace{.15cm}
	
	\only<1>{
		Der Spannungsabfall am letzten Knoten darf nicht zu groß werden.
		\begin{equation*}
	U_0-\sum_{l\in\mathcal{L}}\bigg(Z_l\cdot\sum_{\substack{
	n\in\mathcal{N}\\ n\geq l}}\Big(I^\mathsf{HH}_{t,n}+I^\mathsf{BEV}_{t,n}
	\Big)\bigg)\geq U_\mathsf{min}\quad\forall t\in\mathcal{T}
		\end{equation*}
	}
	
	\only<2>{
		Die Summe der Ströme darf nicht zu groß werden.
		\begin{equation*}
	\sum_{n\in\mathcal{N}}\Big(I^\mathsf{HH}_{t,n}+I^\mathsf{BEV}_{t,n} \Big)
	\leq I_\mathsf{max}\quad\forall t\in\mathcal{T}
		\end{equation*}
	}
	
	\only<3>{
		Die Ströme auf den einzelnen Leitungen dürfen jeweils nicht zu
		groß werden.
		\begin{equation*}
	\sum_{\substack{n\in\mathcal{N}\\n\geq l}}\Big(I^\mathsf{HH}_{t,n}
	+I^\mathsf{BEV}_{t,n} \Big)
	\leq I^\mathsf{max}_l\quad\forall t,l\in\mathcal{T}\times\mathcal{L}
		\end{equation*}
	}
	
	\only<4>{
		Die Energieerhaltung während des Ladevorgangs muss gewährleistet
		sein.
		\begin{equation*}
	SOC_{t,n}+\frac{I^\mathsf{BEV}_{t,n}\cdot U_n\cdot\Delta t}{E_n}
	\cdot 100\%=SOC_{t+1,n}\quad\forall t,n\in\mathcal{T}\times
	\mathcal{N}_\mathsf{BEV}
		\end{equation*}
	}
	
	\only<5>{
		Außerdem vier weitere, optionale Restriktionen für \enquote{faires}
		Laden.
	}
\end{frame}


%\begin{frame}{Rolling Horizon für große Optimierungsprobleme}
%	\begin{itemize}
%		\item Großer Zeithorizont, zu dem ein großes Optimierungsproblem
%		korrespondiert
%		\item Großen Zeithorizont zerlegen in mehrere kleine Horizonte
%		\item Die zu den einzelnen Horizonten korrespondierenden 
%		Optimierungsprobleme lösen
%		\item Die Gesamtlösung ergibt sich aus den einzelnen Teillösungen
%	\end{itemize}
%\end{frame}


\begin{frame}{Rolling Horizon für große Optimierungsprobleme}
	\begin{figure}
		\centering
		%%%%%%%%%%%%%%%%%%%%%%%%%%%%%%%%%%%%%%%%%%%%%%%%%%%%%%%%%%%%%%%%%%%%%%%%%%%%%%%
% Schematische Darstellung eines Rolling Horizon
%%%%%%%%%%%%%%%%%%%%%%%%%%%%%%%%%%%%%%%%%%%%%%%%%%%%%%%%%%%%%%%%%%%%%%%%%%%%%%%

\begin{tikzpicture}
	\draw[->] (0,0) -- (10,0) node[right]{$t$};
	% 1. Horizont
	\node[left] at (0,-0.5) {$\mathcal{T}_1$};
	\foreach \x in {0,1,...,5}{
		\draw(\x,-.75) rectangle (\x+1,-.25);
		\fill[black, opacity={0.6-\x/10}](\x,-.75) rectangle (\x+1,-.25);
	}
	\draw[-stealth] (1.5,-0.5) .. controls (1.8,-.75) and (1.8,-1)
	.. (1.5,-1.25);
		
	% 2. Horizont
	\node[left] at (0,-1.25) {$\mathcal{T}_2$};
	\foreach \x in {1,2,...,6}{
		\draw(\x,-1.5) rectangle (\x+1,-1);
		\fill[black, opacity={0.7-\x/10}](\x,-1.5) rectangle (\x+1,-1);
	}
	\draw[-stealth] (2.5,-1.25) .. controls (2.8,-1.5) and (2.8,-1.75) 
	.. (2.5,-2);
		
	% 3. Horizont
	\node[left] at (0,-2) {$\mathcal{T}_3$};
	\foreach \x in {2,3,...,7}{
		\draw (\x,-2.25) rectangle (\x+1,-1.75);
		\fill[black, opacity={0.8-\x/10}](\x,-2.25) rectangle (\x+1,-1.75);
	}
	\draw[-stealth] (3.5,-2) .. controls (3.8,-2.25) and (3.8,-2.5) 
	.. (3.5,-2.75);
	
	% 4. Horizont
	\node[left] at (0,-2.75) {$\mathcal{T}_4$};
	\foreach \x in {3,4,...,8}{
		\draw (\x,-3) rectangle (\x+1,-2.5);
		\fill[black, opacity={0.9-\x/10}](\x,-3) rectangle (\x+1,-2.5);
	}
	
	% oben geschweifte Klammer, welche die Weite eines Horizonts angibt
	\draw[decorate, decoration=brace] (0, 0.3) -- (6, 0.3)
	node[pos=0.5, above=1mm]{$w$};
		
	% unten das Ergebnis zeichnen
	\node[left] at (0,-4) {$\mathcal{T}$};
	\foreach \x in {0,1,...,9}{
		\draw (\x,-3.75) rectangle (\x+1,-4.25);
	}
	
	\foreach \x in {0,1,...,3}{
		\fill[black, opacity={0.6}] (\x,-3.75) rectangle (\x+1,-4.25);
	}
	
	% Pfeile, die anzeigen, dass die jeweils ersten Ergebnisse von
	% den Horizonten übernommen werden
	\draw[-stealth] (0.5,-0.9) -- (0.5,-3.6);
	\draw[-stealth] (1.5,-1.65) -- (1.5,-3.6);
	\draw[-stealth] (2.5,-2.4) -- (2.5,-3.6);
	\draw[-stealth] (3.5,-3.15) -- (3.5,-3.6);
	
	% diagonale Pünktchen, die andeuten, dass noch mehr Horizonte folgen
	\node at (6.5,-3.2){$\ddots$};
	\node[left] at (-0.2,-3.1){$\vdots$};
	
\end{tikzpicture}
	\end{figure}
\end{frame}


\section{Praktische Umsetzung}

\begin{frame}{Optimierungsprobleme mit Python lösen}
	\begin{itemize}
		\item Python als Programmiersprache
		\item pyomo als Optimierungsframework
		\item glpk als Solver
	\end{itemize}
	
	\onslide<2>{
		\begin{block}{Vorteile von pyomo:}
			\begin{itemize}
				\item Keine Begrenzung der Länge der Zielfunktion
				\item Ermöglicht ein dynamisches Erzeugen des
				Optimierungsmodells
			\end{itemize}
		\end{block}
	}
\end{frame}


\begin{frame}{Aufbau der Software zur Optimierung}
	Mehrere Klassen:
	\begin{itemize}
		\item Eigene Klasse \texttt{GridLineOptimizer} für Netzstrahl und 
		Logik der Optimierung
		\item Eigene Klasse \texttt{BatteryElectricVehicle} für BEVs
		\item Eigene Klasse \texttt{Household} für Haushaltslasten
	\end{itemize}
	
	\onslide<2>{
		\begin{block}{Modularer Aufbau}
			Bessere Wartbarkeit und Portabilität
		\end{block}
	}
\end{frame}


\begin{frame}{Funktionsprinzip der Software zur Optimierung}
	\begin{figure}
		\centering
		%%%%%%%%%%%%%%%%%%%%%%%%%%%%%%%%%%%%%%%%%%%%%%%%%%%%%%%%%%%%%%%%%%%%%%%%%%%%%%%
% flow chart für das Funktionsprinzip der Software
%%%%%%%%%%%%%%%%%%%%%%%%%%%%%%%%%%%%%%%%%%%%%%%%%%%%%%%%%%%%%%%%%%%%%%%%%%%%%%%

\begin{tikzpicture}[
	block/.style={
		font=\tiny,
		draw,
		rectangle,
		minimum width={4cm},
		execute at begin node=\strut,
		}
	]
	
	\node[block](grid){Netztopologie};
	
	\node[block,
		right=0.3 of grid](time){Zeitliche Betrachtung};
		
	\node[block,
		below left=0.5 and 0.15 of time.south](sets){Entsprechende Sets 
		erzeugen};
		
	\node[block,
		below=0.3 of sets](params){Parameter erzeugen};
		
	\node[block,
		below=0.3 of params](vars){Variablen erzeugen};
		
	\node[block,
		below=0.3 of vars](bounds){Ober-/Untergrenzen festlegen};
		
	\node[block,
		below=0.3 of bounds](target){Zielfunktion erzeugen};
		
	\node[block,
	 	below=0.3 of target](constraints){Restriktionen erzeugen};
	 	
	 \node[block,
	 	right=0.9 of params](households){Haushaltslastprofile};
	 	
	 \node[block,
	 	right=0.9 of bounds](cust-wish){Kundenwunsch};
	 	
	 \node[block,
	 	below=0.5 of constraints](result){Optimierungsmodel lösen};
	 	
	 \node[left=0.3 of sets,
	 	rotate=90](model){\scriptsize Optimierungsmodel erzeugen};
	 	
	 \node[above=0.1 of cust-wish](BEV){\footnotesize\ttfamily 
	 BatteryElectricVehicle};
	 
	 \node[above=0.1 of households](HH){\footnotesize\ttfamily Household};
	 
	 \node[left=1 of params,
	 	rotate=90](GLO){\footnotesize\ttfamily GridLineOptimizer};
	 	
	 	
	 \draw[-stealth](grid)--(sets);
	 \draw[-stealth](time)--(sets);
	 \draw[-stealth](sets)--(params);
	 \draw[-stealth](params)--(vars);
	 \draw[-stealth](vars)--(bounds);
	 \draw[-stealth](bounds)--(target);
	 \draw[-stealth](target)--(constraints);
	 \draw[-stealth](constraints)--(result);
	 \draw[-stealth](cust-wish)--(bounds);
	 \draw[-stealth](households)--(params);
	 
	 \begin{pgfonlayer}{background}
	 	\path(sets.west)+(-0.6, 0.5) node (a) {};
		\path(constraints.east)+(0.1, -0.5) node (b) {};
		\path[fill=yellow!20, rounded corners, draw=black!50, dashed]
		(a) rectangle (b);
	 \end{pgfonlayer}

\end{tikzpicture}
	\end{figure}
\end{frame}


\begin{frame}{Schnittstelle zwischen Optimierung und Simulationsumgebung}
	\begin{itemize}
		\item Bisher liefert der Optimierungsalgorithmus nur optimierte
		Ladeströme
		\item Mithilfe der Simulationsumgebung können die Ergebnisse
		der Optimierung validiert werden
		\item Hierzu wird eine Schnittstelle zwischen Optimierung und
		Simulationsumgebung geschaffen
	\end{itemize}
\end{frame}


\begin{frame}{Funktionsprinzip der Schnittstelle}
	\begin{figure}
		\centering
		%%%%%%%%%%%%%%%%%%%%%%%%%%%%%%%%%%%%%%%%%%%%%%%%%%%%%%%%%%%%%%%%%%%%%%%%%%%%%%%
% Schnttstelle GLO zu EMO
%%%%%%%%%%%%%%%%%%%%%%%%%%%%%%%%%%%%%%%%%%%%%%%%%%%%%%%%%%%%%%%%%%%%%%%%%%%%%%%

\begin{tikzpicture}[
	block/.style={
		font=\scriptsize,
		draw,
		rectangle,
		minimum width={3.5cm},
		align=center,
		}
	]
	
	\node[block,
		minimum height=6cm](GLO){Optimierer\\ \texttt{GridLineOptimizer}};
		
	\node[block,
		right=6 of GLO,
		minimum height=6cm](EMO){Simulationsumgebung\\ 
		\texttt{Simulation\_Handler}};
		
	\draw[-stealth](GLO.55) -- node[auto] {\scriptsize Netztopologie 
	exportiere} node[below]{}(EMO.125);
		
	\draw[-stealth](GLO.35) -- node[auto] {\scriptsize Haushalts-Profile 
	exportieren} node[below]{} (EMO.145);
		
	\draw[-stealth](GLO.0) -- node[auto] {\scriptsize Ergebnis der Ströme
	exportieren} node[below]{} (EMO.180);
		
	\draw[-stealth](EMO.215) -- node[above] {\scriptsize Ergebnis der 
	Knotenspannungen exportieren} node[below]{} (GLO.325);

\end{tikzpicture}

	\end{figure}
\end{frame}


\begin{frame}{Paralellisierung von Optimierung und Simulationsumgebung}
	\begin{itemize}
		\item Bisher ist nur eine serielle Ausführung von Optimierung und
		Simulationsumgebung möglich
		\item Solange die Optimierung ausgeführt wird, kann nicht reagiert 
		werden
		\item Deshalb wird eine parallele Ausführung von Optimierung und
		Simulationsumgebung realisiert
	\end{itemize}
\end{frame}


\begin{frame}{Paralellisierung von Optimierung und Simulationsumgebung}
	\begin{figure}
		\centering
		%%%%%%%%%%%%%%%%%%%%%%%%%%%%%%%%%%%%%%%%%%%%%%%%%%%%%%%%%%%%%%%%%%%%%%%%%%%%%%%
% Prinzip der Parallelisierung
%%%%%%%%%%%%%%%%%%%%%%%%%%%%%%%%%%%%%%%%%%%%%%%%%%%%%%%%%%%%%%%%%%%%%%%%%%%%%%%

\begin{tikzpicture}[
		block/.style={
			font=\small,
			draw,
			rectangle,
			minimum width={2.0cm},
			minimum height={1.2cm},
			align=center
			}
		]
		
		\node foreach \t in {1,3,5,7,9,11} 
		[draw, block](o\t) at (\t,1.0){O};
		
		\node foreach \t in {1,...,24} 
		[draw, block,
		minimum width={0.5cm}](u\t) at ({\t*0.5-0.25},-1.0){S};
		
		\draw[-stealth](0,0)--(13,0) node[right] {\small t};
		\draw[->](o1.south east)--(u4.north east);
		\draw[->](o3.south east)--(u8.north east);
		\draw[->](o5.south east)--(u12.north east);
		\draw[->](o7.south east)--(u16.north east);
		\draw[->](o9.south east)--(u20.north east);
		
		\draw (u1.south) edge[out=270, in=270, ->] (u2.south);
		\draw (u2.south) edge[out=270, in=270, ->] (u3.south);
		\draw (u3.south) edge[out=270, in=270, ->] (u4.south);
		
		\draw (u5.south) edge[out=270, in=270, ->] (u6.south);
		\draw (u6.south) edge[out=270, in=270, ->] (u7.south);
		\draw (u7.south) edge[out=270, in=270, ->] (u8.south);
		
		\node[below=0.01 of u9](dots){\tiny \dots};
		
		
		\draw[-stealth](0,4.5)--(13,4.5) node[right] {\small t};
		
		\node foreach \t in {1,3.5,6,8.5,11}
		[draw, block] at (\t,3.5){O};
		
		\node foreach \t in {2.25,4.75,7.25,9.75}
		[draw, block,
		minimum width={0.5cm}] at (\t,3.5){S};
		
	\end{tikzpicture}
	\end{figure}
\end{frame}


\begin{frame}{Prinzip der Paralellisierung}
	\begin{figure}
		\centering
		%%%%%%%%%%%%%%%%%%%%%%%%%%%%%%%%%%%%%%%%%%%%%%%%%%%%%%%%%%%%%%%%%%%%%%%%%%%%%%%
% Ablauf vom Code der Parallelisierung
%%%%%%%%%%%%%%%%%%%%%%%%%%%%%%%%%%%%%%%%%%%%%%%%%%%%%%%%%%%%%%%%%%%%%%%%%%%%%%%

\begin{tikzpicture}[
	block/.style={
		font=\tiny,
		draw,
		rectangle,
		text width={4.0cm},
		align=center,
		execute at begin node=\strut
		}
	]
	
	\node[block](GLO){Objekte zur Optimierung erzeugen};	
	
	\node[block,
		below=0.4 of GLO](exportGrid){Zur Optimierung verwendetes Netz 
		exportieren};
	
	\node[block,
		below=0.4 of exportGrid](update){Indizes sowie Ober- und Untergrenzen 
		anpassen};
	
	\node[block,
		below=0.4 of update](build){Optimierungsmodell entsprechend aufbauen};
		
	\node[block,
		below=0.4 of build](exportResults){Aktuelle Ergebnisse aufbereiten};
		
	\node[block,
		right=0.8 of exportResults,
		text width=2cm](queue){Queue};
		
	\node[block,
		right=0.8 of queue](receiveResults){Aktuelles Ergebnis aus Queue 
		auslesen};
		
	\node[block,
		below=0.4 of receiveResults](simulate){Simulation mit aktuellen 
		Werten};
		
	\node[block] (receiveGrid) at (receiveResults |- exportGrid)
		{Netz gemäß den Angaben rekonstruieren};
		
	\node[below=0.1 of exportResults](optProc){\tiny Prozess der 
	Optimierung};
	
	\node[below=0.1 of simulate](simProc){\tiny Prozess der 
	Netzsimulation};
		
	
	\draw[-stealth](GLO) -- (exportGrid);	
	\draw[-stealth](exportGrid)--(receiveGrid);
	\draw[-stealth](exportResults)--(queue);
	\draw[-stealth](queue)--(receiveResults);
	\draw[-stealth](exportGrid)--(update);
	\draw[-stealth](update)--(build);
	\draw[-stealth](build)--(exportResults);
	\draw[-stealth](exportResults.west)-- ++(-0.5,0) |- (update.west);
	\draw[-stealth](simulate.east)-- ++(0.5,0) |- (receiveResults.east);
	\draw[-stealth](receiveResults) -- (simulate);
	\draw[-stealth](receiveGrid) -- (receiveResults);
	
	
	\begin{pgfonlayer}{background}
		\path(update.north west)+(-0.8, 0.1) node (a) {};
		\path(exportResults.east)+(0.1, -0.8) node (b) {};
		\path[fill=yellow!20, rounded corners, draw=black!50, dashed]
		(a) rectangle (b);
		
		\path(simulate.east)+(0.8, -0.8) node (a) {};
		\path(receiveResults.north west)+(-0.1, 0.1) node (b) {};
		\path[fill=yellow!20, rounded corners, draw=black!50, dashed]
		(a) rectangle (b);
	\end{pgfonlayer}
	
\end{tikzpicture}
	\end{figure}
\end{frame}


\section{Ergebnisse}

\begin{frame}{Untersuchung anhand von Szenarien}
	\begin{itemize}
		\item Definieren unterschiedlicher Szenarien
		\begin{itemize}
			\item Transformatorleistung
			\item Anzahl Knoten
			\item Position der Ladesäulen
			\item Kundenwünsche
			\item $\dots$
		\end{itemize}
		\item Optimierung ausführen
		\item Ergebnisse mithilfe der Simulationsumgebung validieren
	\end{itemize}
\end{frame}


\begin{frame}{Festlegen von Bewertungskriterien}
	\begin{itemize}
		\item Kriterien der Netzstabilität
		\begin{itemize}
			\item Auslastung des Transformators
			\item Auslastung der Leitungen
			\item Spannung am letzten Knoten
		\end{itemize}
		\item Kundenwünsche erfüllen
		\item \enquote{Faires} Laden ermöglichen
	\end{itemize}
\end{frame}


\section{Fazit}

\begin{frame}{Fazit}
	Punkt 1 der Aufgabenstellung:
	\begin{itemize}
		\item Skalierbaren Optimierungsalgorithmus entwickelt
		\item Ausführbarkeit auf Raspberry Pi demonstriert
		\item Optimierung allgemeiner Netztopologien muss noch umgesetzt werden
	\end{itemize}\vspace{.5cm}
	Punkt 2 der Aufgabenstellung:
	\begin{itemize}
		\item Schnittstelle zwischen Optimierung und Simulationsumgebung
		geschaffen
		\item Sinnhaftigkeit der Ergebnisse der Optimierung validiert
	\end{itemize}\vspace{.5cm}
	Punkt 3 der Aufgabenstellung:
	\begin{itemize}
		\item Parallelisierung von Optimierung und Netzsimulation realisiert
		\item bla
	\end{itemize}
\end{frame}


\appendix
\setbeamertemplate{footline}[default]
\begin{withoutheadline}
\begin{frame}[noframenumbering]{Literatur}
	\printbibliography[heading=none]
\end{frame}
\end{withoutheadline}


\begin{withoutheadline}
\begin{frame}[noframenumbering]{Vielen Dank für die Aufmerksamkeit!}
	\begin{columns}
		\begin{column}{.3\textwidth}
			Noch Fragen?
		\end{column}
		\begin{column}{.7\textwidth}
			%%%%%%%%%%%%%%%%%%%%%%%%%%%%%%%%%%%%%%%%%%%%%%%%%%%%%%%%%%%%%%%%%%%%%%%%%%%%%%%
% flow chart für das Funktionsprinzip der Software
%%%%%%%%%%%%%%%%%%%%%%%%%%%%%%%%%%%%%%%%%%%%%%%%%%%%%%%%%%%%%%%%%%%%%%%%%%%%%%%

\begin{tikzpicture}[
	block/.style={
		font=\tiny,
		color=white,
		align=left
		}
	]
	\node (monkey) {%
		\includegraphics[width=\textwidth]{Affe}
	};
	
	\node[block, below of=monkey, yshift=-2cm]
	{\url{https://imgflip.com/memetemplate/28186299/monkey-scratch}};
\end{tikzpicture}
		\end{column}
	\end{columns}
\end{frame}
\end{withoutheadline}


\begin{withoutheadline}
\begin{frame}[noframenumbering]{Formulierung als Linear Programming}
	Die Zielfunktion: Maximieren der Ladeströme zu allen Zeitpunkten
	\begin{equation*}
		\text{max.}\quad\sum_{t\in\mathcal{T}}\bigg(\sum_{n\in
		\mathcal{N_\mathsf{BEV}}}\!\!I^\mathsf{BEV}_{t,n}\bigg)
	\end{equation*}
\end{frame}
\end{withoutheadline}

\end{document}